% Companion correction section. General parity is known; these explicit
% rows corroborate corrections to the pinned source snapshot.

\section{The proposed mixed-root survivors reduce exactly}
\label{sec:mixed-parity-correction}

The snapshot identifies two purportedly new mixed directions in each of
the Gaussian and Eisenstein towers: a real depth-two value of weight five
and an imaginary depth-two value of weight six. Both are in the component
covered by the established parity theorem. Their resistance to a numerical
integer-relation search is therefore not evidence of a new direction in
the algebra of single cyclotomic values. Explicit formulas settle all four
cases. These formulas are presented as corroborated corrections to the
snapshot. The general parity theorem predates this work, and recent
companion work has already reported four mixed-root reductions; no claim
of first discovery is made here.

For $x=e^{i\theta}$, $0<\theta<2\pi$, put
\[
 M_{a,b}(x)=\operatorname{Li}_{a,b}(x,x^{-1}),\qquad
 \mathcal B_n(x)=\frac{(2\pi i)^n}{n!}B_n\!\left(\frac{\theta}{2\pi}\right),
 \qquad
 \kappa_0=1,\quad\kappa_j=-2\zeta(2j)\ (j\ge1).
\]
Use $C_{\mu,r}=\binom{\mu-1}{r-1}$ when $\mu\ge r$ and zero otherwise.

\begin{theorem}[Inverse-color parity formula]
\label{thm:inverse-color}
For all $a,b\ge1$ and $w=a+b$,
\begin{align}
 M_{a,b}(x)-(-1)^w\overline{M_{a,b}(x)}
 ={}&(-1)^a\sum_{j=0}^{\lfloor(w-1)/2\rfloor}
 \kappa_j\bigl(C_{w-2j,a}+C_{w-2j,b}\bigr)
 \operatorname{Li}_{w-2j}(\bar x)\label{eq:inverse-color}\\
 &-\zeta(w)-\operatorname{Li}_b(\bar x)\mathcal B_a(x).\notag
\end{align}
Thus the real part at odd weight and the imaginary part at even weight
reduce to single polylogarithms times zeta values and Bernoulli factors.
\end{theorem}

\begin{proof}
In Panzer's depth-two inversion formula, use the ascending-index pair
$(b,a)$ and arguments $(y,x)$, and let $y\to x^{-1}$. Only the
Bernoulli factor at $xy=1$ has a possible side-dependent value. Its
index-one coefficient is
\[
 \left[(-1)^{b+w-1}\binom{w-2}{b-1}
       +(-1)^a\binom{w-2}{a-1}\right]
       \operatorname{Li}_{w-1}(x^{-1})=0.
\]
The binomial coefficients agree and the two signs are opposite.
All odd Bernoulli numbers of index greater than one vanish; the index-zero
value is $1$ and the even positive values are $-2\zeta(2j)$.
In the surviving terms $\mu=w-2j$, the first sum's sign
$(-1)^{b+\mu}$ equals $(-1)^a$. The two sums combine to
\eqref{eq:inverse-color}. The single polylogarithm at the product tends to
$\zeta(w)$, since $w\ge2$. The remaining terms are analytic at $y=x^{-1}$,
because $x\ne1$. The defining boundary double sums also have the stated
limit: when $a>1$ they converge absolutely, while for $a=1$ the bounded
inner partial sums are a constant plus $O(n^{-b})$, leaving a convergent
Dirichlet series and an absolutely convergent remainder. Complex
conjugation now identifies the inverse pair with
$\overline{M_{a,b}(x)}$.
\end{proof}

At $b=1$ the formula simplifies substantially:
\begin{align}
 M_{a,1}(x)-(-1)^{a+1}\overline{M_{a,1}(x)}
 ={}&(-1)^a\left[(a+1)\operatorname{Li}_{a+1}(\bar x)
       -2\sum_{j=1}^{\lfloor a/2\rfloor}
          \zeta(2j)\operatorname{Li}_{a+1-2j}(\bar x)\right]
 \label{eq:inverse-color-one}\\
 &-\zeta(a+1)-\operatorname{Li}_1(\bar x)\mathcal B_a(x).\notag
\end{align}
This single equation proves all four corrections below.

\subsection{Gaussian corrections}

Using $\mathcal B_4(i)=7\pi^4/5760$,
$\mathcal B_5(i)=-5i\pi^5/768$, and
$\operatorname{Li}_1(-i)=-\log2/2-i\pi/4$ gives
\begin{align}
 \Re\operatorname{Li}_{4,1}(i,-i)
  &=\frac{3\pi^4\log2}{512}+\frac{\pi^2\zeta(3)}{64}
                          -\frac{587\zeta(5)}{1024},
 \label{eq:gaussian-mixed-real}\\
 \Im\operatorname{Li}_{5,1}(i,-i)
  &=3\beta(6)-\frac{\pi^2\beta(4)}6-\frac{\pi^4G}{90}
                          -\frac{5\pi^5\log2}{3072}.
 \label{eq:gaussian-mixed-imag}
\end{align}
Every value on the right belongs to the single-value basket already named
in the source's Gaussian search. Consequently these two rows cannot add
new directions modulo that basket.

\subsection{Eisenstein corrections}

Let $\rho=e^{2\pi i/3}$ and write
$C_j=\Im\operatorname{Li}_j(\rho)$ for even $j$.
Here $x=\rho^2$, so the Bernoulli argument is $2/3$:
\[
 \mathcal B_4(\rho^2)=\frac{13\pi^4}{1215},\qquad
 \mathcal B_5(\rho^2)=\frac{4i\pi^5}{729},\qquad
 \operatorname{Li}_1(\rho)=-\frac12\log3+\frac{i\pi}{6}.
\]
Equation~\eqref{eq:inverse-color-one} gives
\begin{align}
 \Re\operatorname{Li}_{4,1}(\rho^2,\rho)
  &=\frac{2\pi^4\log3}{243}+\frac{2\pi^2\zeta(3)}{27}
                           -\frac{281\zeta(5)}{162},
 \label{eq:eisenstein-mixed-real}\\
 \Im\operatorname{Li}_{5,1}(\rho^2,\rho)
  &=-3C_6+\frac{\pi^2C_4}{6}+\frac{\pi^4C_2}{90}
                         +\frac{\pi^5\log3}{729}.
 \label{eq:eisenstein-mixed-imag}
\end{align}
The corresponding imaginary formula at $(\rho,\rho^2)$ has the opposite
sign by conjugation. Again all right-hand atoms already occur in the
source's stated search basket.

\paragraph{Precise editorial consequence.}
The claimed parallel location of candidate new mixed directions does not
survive these exact reductions. The text should say that the tested
convergent double-shuffle system left formal residual coordinates which
are eliminated by known parity relations. Numerical or formal residual
counts for larger systems must be recomputed after adjoining these rows;
one should not infer a corrected total numerical dimension merely by
subtracting two from a previously asserted count. The established
inversion theorem supplies reductions, not independence statements for
the opposite parity components.

The real weight-five rows give a particularly concrete audit target:
\[
 1024\Re\operatorname{Li}_{4,1}(i,-i)
   =6\pi^4\log2+16\pi^2\zeta(3)-587\zeta(5),
\]
\[
 486\Re\operatorname{Li}_{4,1}(\rho^2,\rho)
   =4\pi^4\log3+36\pi^2\zeta(3)-843\zeta(5).
\]
Their integer-relation heights, including the target coefficient, are
$1024$ and $843$. Both are below the source's stated approximate
$10^4$ search scale. Missing these relations therefore cannot be
attributed solely to very large coefficients. The actual input values,
basis construction, dependence handling, and stopping rules should be
reproduced; without the original scripts, it would be premature to
identify one particular cause.

\paragraph{Independent numerical check.}
The exact formulas in \texttt{mixed\_candidates.py} were also compared
with the independent radial integral
\[
 \operatorname{Li}_{a,b}(x,x^{-1})
 =\frac{x}{(a-1)!}\int_0^1
       \frac{(-\log t)^{a-1}\operatorname{Li}_b(t)}{1-xt}\,dt.
\]
At 60 decimal digits all four absolute residuals were below $10^{-60}$.
For orientation, the four left sides in the order above are approximately
$-0.0134202673573$, $0.0326942444522$,
$-0.0390392184565$, and $-0.0253393813681$.
The analytic formula proves the identities; these numerical checks serve
only to detect arithmetic or convention errors.

